\chapter{Análisis y diseño}
\label{chap:analisis_diseno}

\lettrine{E}{n} este capítulo se presenta el análisis de requisitos del framework desarrollado, la arquitectura general del sistema y el diseño detallado de cada uno de los módulos componentes. El diseño se basa en una metodología incremental que ha evolucionado desde prototipos iniciales hasta una arquitectura modular robusta que cumple con los principios \acrshort{solid} y patrones de diseño modernos.

\section{Análisis de requisitos}
\label{sec:analisis_requisitos}

\subsection{Requisitos funcionales}
\label{subsec:requisitos_funcionales}

Los requisitos funcionales del sistema se organizan en cuatro categorías principales, correspondientes a los módulos del framework:

\subsubsection{Captura de eventos de entrada}
\label{subsubsec:req_captura_entrada}

\begin{enumerate}
    \item \textbf{RF-01}: El sistema debe capturar eventos de teclado con precisión temporal de milisegundos
    \item \textbf{RF-02}: El sistema debe capturar eventos de ratón (botones, movimiento, rueda) con múltiples modos de seguimiento
    \item \textbf{RF-03}: El sistema debe soportar controladores PS4 con captura de botones, sticks analógicos y triggers
    \item \textbf{RF-04}: El sistema debe proporcionar visualización en tiempo real de todos los dispositivos de entrada
    \item \textbf{RF-05}: El sistema debe permitir configuración dinámica de teclas/botones a monitorizar
    \item \textbf{RF-06}: El sistema debe registrar eventos en múltiples formatos (CSV, JSON) simultáneamente
\end{enumerate}

\subsubsection{Procesamiento de señales EEG}
\label{subsubsec:req_procesamiento_eeg}

\begin{enumerate}
    \item \textbf{RF-07}: El sistema debe soportar múltiples dispositivos EEG comerciales mediante BrainFlow
    \item \textbf{RF-08}: El sistema debe implementar filtros digitales en tiempo real (pasa-banda, notch, detrend)
    \item \textbf{RF-09}: El sistema debe proporcionar análisis avanzado mediante \acrshort{ica}, transformadas wavelet y downsampling
    \item \textbf{RF-10}: El sistema debe integrar MNE-Python para análisis neurocientífico avanzado
    \item \textbf{RF-11}: El sistema debe visualizar señales EEG en tiempo real con mapas topográficos cerebrales
    \item \textbf{RF-12}: El sistema debe evaluar calidad de señal con indicadores visuales
    \item \textbf{RF-13}: El sistema debe permitir inserción de marcadores temporales durante la grabación
\end{enumerate}

\subsubsection{Sincronización y exportación de datos}
\label{subsubsec:req_sincronizacion}

\begin{enumerate}
    \item \textbf{RF-14}: El sistema debe sincronizar temporalmente señales EEG y eventos de entrada
    \item \textbf{RF-15}: El sistema debe exportar datos a formatos estándar (\acrshort{bids}, \acrshort{edf}+, \acrshort{hdf5}, \acrshort{csv})
    \item \textbf{RF-16}: El sistema debe utilizar \acrshort{lsl} para sincronización de alta precisión
    \item \textbf{RF-17}: El sistema debe generar metadatos completos para datasets exportados
    \item \textbf{RF-18}: El sistema debe proporcionar interfaz unificada para coordinar ambos módulos
\end{enumerate}

\subsubsection{APIs y interfaces de usuario}
\label{subsubsec:req_interfaces}

\begin{enumerate}
    \item \textbf{RF-19}: El sistema debe proporcionar interfaz gráfica responsiva basada en PyQt5
    \item \textbf{RF-20}: El sistema debe implementar \acrshort{cli} con capacidades de monitoreo en tiempo real
    \item \textbf{RF-21}: El sistema debe exponer \acrshort{api} \acrshort{rest} para control remoto
    \item \textbf{RF-22}: El sistema debe proporcionar WebSocket para streaming de datos en tiempo real
    \item \textbf{RF-23}: El sistema debe incluir explorador web de APIs con documentación interactiva
    \item \textbf{RF-24}: El sistema debe soportar modo headless para uso programático
\end{enumerate}

\subsection{Requisitos no funcionales}
\label{subsec:requisitos_no_funcionales}

\subsubsection{Rendimiento}
\label{subsubsec:req_rendimiento}

\begin{itemize}
    \item \textbf{RNF-01}: Latencia de captura de eventos < 50ms
    \item \textbf{RNF-02}: Procesamiento EEG en tiempo real para frecuencias de muestreo hasta 1000 Hz
    \item \textbf{RNF-03}: Sincronización temporal con precisión < 10ms
    \item \textbf{RNF-04}: Soporte para streaming WebSocket con latencia < 15ms
    \item \textbf{RNF-05}: Uso de memoria optimizado con buffers configurables
\end{itemize}

\subsubsection{Usabilidad}
\label{subsubsec:req_usabilidad}

\begin{itemize}
    \item \textbf{RNF-06}: Interfaz gráfica intuitiva con tema oscuro/claro
    \item \textbf{RNF-07}: Documentación contextual integrada (25+ temas de ayuda)
    \item \textbf{RNF-08}: Configuración sin necesidad de programación para usuarios básicos
    \item \textbf{RNF-09}: Retroalimentación visual inmediata del estado del sistema
\end{itemize}

\subsubsection{Extensibilidad}
\label{subsubsec:req_extensibilidad}

\begin{itemize}
    \item \textbf{RNF-10}: Arquitectura modular que permita adición de nuevos colectores
    \item \textbf{RNF-11}: Interfaces bien definidas para integración con herramientas externas
    \item \textbf{RNF-12}: Sistema de plugins para algoritmos de procesamiento personalizados
    \item \textbf{RNF-13}: Compatibilidad con formatos de datos emergentes
\end{itemize}

\subsubsection{Calidad}
\label{subsubsec:req_calidad}

\begin{itemize}
    \item \textbf{RNF-14}: Código con tipado estático completo (mypy compliance)
    \item \textbf{RNF-15}: Documentación de código con docstrings (pydocstyle compliance)
    \item \textbf{RNF-16}: Formateo automático de código (black compliance)
    \item \textbf{RNF-17}: Thread-safety en todas las operaciones concurrentes
\end{itemize}

\section{Arquitectura general del sistema}
\label{sec:arquitectura_general}

\subsection{Visión general}
\label{subsec:vision_general}

El framework implementa una \textbf{arquitectura orientada a servicios} con \textbf{comunicación dirigida por eventos}, organizada en cuatro módulos principales que pueden operar independientemente o de forma coordinada.

\subsection{Principios arquitectónicos}
\label{subsec:principios_arquitectonicos}

\subsubsection{Separación de responsabilidades}
Cada módulo tiene responsabilidades claramente definidas:

\begin{itemize}
    \item \textbf{input\_collector}: Captura y procesamiento de eventos de entrada
    \item \textbf{eeg\_processor}: Adquisición y análisis de señales EEG
    \item \textbf{api\_explorer}: Framework web para exploración de APIs
    \item \textbf{synchronization\_module}: Coordinación temporal y exportación de datasets
\end{itemize}

\subsubsection{Comunicación desacoplada}
Los módulos se comunican mediante:

\begin{itemize}
    \item \textbf{Core Event System}: Sistema de eventos independiente de framework
    \item \textbf{APIs REST}: Interfaces well-defined para comunicación inter-módulo
    \item \textbf{WebSocket}: Streaming de datos en tiempo real
    \item \textbf{LSL Streams}: Sincronización temporal de alta precisión
\end{itemize}

\subsubsection{Arquitectura multi-interfaz}
Cada módulo soporta múltiples interfaces simultáneamente:

\begin{itemize}
    \item \textbf{GUI}: Interfaces PyQt5 con diseño responsivo
    \item \textbf{CLI}: Interfaces Typer con Rich para output avanzado
    \item \textbf{API REST}: Control programático y integración externa
    \item \textbf{WebSocket}: Monitoreo en tiempo real
\end{itemize}

\subsection{Patrones arquitectónicos implementados}
\label{subsec:patrones_arquitectonicos_impl}

\subsubsection{Event-Driven Architecture}
Sistema de eventos centralizado que facilita:

\begin{itemize}
    \item \textbf{Desacoplamiento}: Componentes independientes
    \item \textbf{Extensibilidad}: Fácil adición de suscriptores
    \item \textbf{Thread-safety}: Comunicación segura entre hilos
    \item \textbf{Framework-agnostic}: Independiente de Qt u otros frameworks
\end{itemize}

\subsubsection{Singleton Pattern}
Servicios centralizados con:

\begin{itemize}
    \item \textbf{Thread-safe implementation}: Double-checked locking
    \item \textbf{Lazy initialization}: Inicialización bajo demanda
    \item \textbf{Resource management}: Gestión centralizada de recursos
\end{itemize}

\subsubsection{Model-View-Controller (MVC)}
Implementado en interfaces complejas:

\begin{itemize}
    \item \textbf{Separación de lógica}: Modelo, vista y control independientes
    \item \textbf{Reutilización}: Modelos compartidos entre vistas
    \item \textbf{Testabilidad}: Componentes aislados para pruebas
\end{itemize}

\subsubsection{Observer Pattern}
Para actualizaciones reactivas:

\begin{itemize}
    \item \textbf{Qt Signals/Slots}: Comunicación UI thread-safe
    \item \textbf{Callback systems}: Notificaciones de cambios de estado
    \item \textbf{Event broadcasting}: Distribución de eventos a múltiples consumidores
\end{itemize}

\section{Diseño del módulo api\_explorer}
\label{sec:diseno_api_explorer}

\subsection{Propósito y responsabilidades}
\label{subsec:proposito_api_explorer}

API Explorer es un framework web modular diseñado para proporcionar capacidades dinámicas de exploración y prueba de APIs para módulos Python. Construido sobre FastAPI con soporte WebSocket, ofrece una interfaz moderna y responsiva que permite a múltiples módulos registrar sus APIs en tiempo de ejecución, proporcionando formularios de parámetros interactivos, transmisión de eventos en tiempo real y visualización de datos. El sistema implementa un patrón singleton para accesibilidad global manteniendo operaciones thread-safe y una arquitectura escalable.

\subsection{Arquitectura del módulo}
\label{subsec:arquitectura_api_explorer}

El API Explorer sigue una \textbf{arquitectura de servicio basada en singleton} combinada con patrones \textbf{observer} y \textbf{registry}, organizando la funcionalidad en capas bien definidas como se muestra en la figura~\ref{fig:arquitectura_api_explorer}.

\begin{figure}[hp!]
\centering
\includegraphics[width=0.9\textwidth,height=15cm,keepaspectratio]{imaxes/api_explorer_capas.pdf}
\caption{Arquitectura en capas del módulo api\_explorer}
\label{fig:arquitectura_api_explorer}
\end{figure}

La arquitectura se estructura en cinco capas principales:

\subsubsection{Capa Cliente}
Múltiples tipos de clientes pueden interactuar con el sistema:
\begin{itemize}
    \item \textbf{Navegador Web Modular}: Frontend JavaScript basado en managers
    \item \textbf{Cliente HTTP (REST)}: Acceso programático a APIs
    \item \textbf{Cliente WebSocket}: Comunicación en tiempo real
\end{itemize}

\subsubsection{Capa Presentación}
Gestiona la interfaz de usuario y experiencia:
\begin{itemize}
    \item \textbf{Plantillas HTML (Jinja2)}: Generación dinámica de contenido
    \item \textbf{CSS Responsivo}: Sistema de temas dual (oscuro/claro)
    \item \textbf{Módulos JavaScript}: Gestión modular de funcionalidad frontend
\end{itemize}

\subsubsection{Framework Web}
Infraestructura de servidor y comunicación:
\begin{itemize}
    \item \textbf{Servidor FastAPI (uvicorn)}: Servidor web asíncrono
    \item \textbf{WebSocket Manager (Worker Pool)}: Gestión distribuida de conexiones en tiempo real
    \item \textbf{Servicio de Archivos Estáticos}: Distribución de recursos frontend
\end{itemize}

\subsubsection{Servicio Core}
Núcleo de la funcionalidad del framework:
\begin{itemize}
    \item \textbf{ApiExplorer (Singleton)}: Coordinador central del sistema
    \item \textbf{Registro de Módulos}: Sistema de registro dinámico
    \item \textbf{Difusión de Eventos}: Distribución de eventos a múltiples clientes
    \item \textbf{Serialización JSON}: Optimización con orjson cuando disponible
    \item \textbf{Validación de Parámetros}: Validación básica de entrada
    \item \textbf{Gestión de Conexiones}: Administración del ciclo de vida de conexiones
\end{itemize}

\subsubsection{Capa Datos}
Estructuras de datos en memoria:
\begin{itemize}
    \item \textbf{Diccionarios de Módulos}: Metadatos de módulos registrados
    \item \textbf{Diccionarios de Métodos}: Definiciones de métodos API
    \item \textbf{Diccionarios de Parámetros}: Especificaciones de parámetros
\end{itemize}

\subsection{Componentes core y su diseño}
\label{subsec:componentes_core_api_explorer}

\subsubsection{API Explorer Core}

El componente central \textbf{ApiExplorer} implementa un patrón singleton thread-safe que coordina todas las operaciones del framework:

\textbf{Características principales:}
\begin{itemize}
    \item \textbf{Patrón Singleton}: Instancia única thread-safe para acceso global
    \item \textbf{Registro Dinámico}: Registro de módulos y métodos en tiempo de ejecución
    \item \textbf{Operaciones Thread-Safe}: Lock principal para sincronización de operaciones críticas
    \item \textbf{Gestión de Puertos}: Detección automática de puertos disponibles
    \item \textbf{Ciclo de Vida del Servidor}: Gestión completa de inicio/parada del servidor
\end{itemize}

El servicio mantiene una estructura de datos jerárquica que organiza módulos, grupos de métodos y métodos individuales, permitiendo un registro flexible y una navegación intuitiva desde la interfaz web.

\subsubsection{Sistema WebSocket con Worker Pools}

El \textbf{WebSocketManager} implementa una arquitectura distribuida para manejar conexiones de alta frecuencia, especialmente crítica para datos EEG en tiempo real:

\textbf{Componentes del sistema:}
\begin{itemize}
    \item \textbf{WebSocketManager}: Coordinador principal de conexiones y workers
    \item \textbf{WebSocketWorkerPool}: Pool de workers para procesamiento distribuido
    \item \textbf{WorkerStats}: Métricas de rendimiento por worker
    \item \textbf{FastJSONSerializer}: Serialización JSON optimizada
\end{itemize}

\textbf{Características de rendimiento:}
\begin{itemize}
    \item \textbf{Arquitectura Worker Pool}: Distribución de carga entre múltiples workers
    \item \textbf{Throttling Adaptativo}: Limitación de mensajes basada en frecuencia
    \item \textbf{Load Balancing}: Distribución round-robin entre workers
    \item \textbf{Performance Tracking}: Métricas de procesamiento en tiempo real
\end{itemize}

\subsubsection{Serialización JSON Optimizada}

El sistema de serialización implementa optimizaciones específicas para el manejo de datos de alta frecuencia:

\textbf{Optimizaciones implementadas:}
\begin{itemize}
    \item \textbf{Integración orjson}: Uso de orjson cuando está disponible para velocidad 2-5x superior
    \item \textbf{Template Caching}: Pre-compilación de estructuras JSON frecuentes
    \item \textbf{Fallback Safety}: Degradación grácil a stdlib json si orjson no está disponible
    \item \textbf{Data Sanitization}: Limpieza automática de datos no serializables
\end{itemize}

\subsection{Arquitectura frontend modular}
\label{subsec:arquitectura_frontend_api_explorer}

\subsubsection{Sistema de managers JavaScript}

El frontend implementa una arquitectura modular basada en managers especializados que se comunican mediante eventos:

\textbf{Managers principales:}
\begin{itemize}
    \item \textbf{APIExplorer Coordinator}: Coordinación central de todos los managers
    \item \textbf{DarkModeManager}: Gestión centralizada de temas oscuro/claro con persistencia
    \item \textbf{WebSocketManager}: Gestión de múltiples conexiones WebSocket con buffering
    \item \textbf{ChartManager}: Coordinación de múltiples instancias Chart.js con updates por lotes
    \item \textbf{EventDisplayManager}: Display de eventos en tiempo real con auto-scrolling
    \item \textbf{FileBrowserManager}: Integración con navegación de archivos del servidor
    \item \textbf{APIMethodsManager}: Ejecución de métodos API y gestión de parámetros dinámicos, incluyendo validación avanzada y carga automática de plantillas JSON
\end{itemize}

\subsubsection{Comunicación event-driven}

Los managers se comunican mediante un sistema de eventos que facilita el desacoplamiento y la extensibilidad. Por ejemplo, cambios en el modo oscuro se propagan automáticamente a los gráficos, y datos recibidos por WebSocket se distribuyen automáticamente a visualizadores y displays de eventos.

\subsection{Visualización de datos y características avanzadas}
\label{subsec:visualizacion_api_explorer}

\subsubsection{Sistema de visualización}

El módulo implementa un sistema completo de visualización de datos con:

\textbf{Integración de gráficos:}
\begin{itemize}
    \item \textbf{Framework Chart.js}: Gráficos profesionales con actualizaciones en tiempo real
    \item \textbf{Streaming en Tiempo Real}: Actualizaciones de datos en vivo con throttling inteligente
    \item \textbf{Soporte Multi-canal}: Visualización separada o combinada para diferentes canales
    \item \textbf{Batched Updates}: Actualizaciones por lotes para suavidad visual
\end{itemize}

\textbf{Modos de visualización:}
\begin{itemize}
    \item \textbf{Gráficos Combinados}: Todos los canales en un solo gráfico
    \item \textbf{Gráficos Separados}: Gráficos individuales por canal de datos
    \item \textbf{Filtrado por Tipo}: Filtrar display por tipo de datos (EEG, EMG, eventos)
    \item \textbf{Auto-escalado}: Escalado dinámico basado en rangos de datos
\end{itemize}

\subsubsection{Formularios de parámetros dinámicos}

El sistema genera automáticamente formularios web basados en las especificaciones de parámetros de cada método API:

\textbf{Tipos de parámetros soportados:}
\begin{itemize}
    \item \textbf{Texto}: Inputs de texto básicos con validación
    \item \textbf{Números}: Inputs numéricos con validación y conversión automática
    \item \textbf{Rangos}: Sliders con formateo y validación de rango
    \item \textbf{Selección}: Dropdowns con opciones predefinidas y disparo de carga de plantillas
    \item \textbf{Archivos}: Navegador de archivos integrado con el servidor
    \item \textbf{JSON Templates}: Áreas de texto con carga automática de plantillas, validación y formateo interactivo
    \item \textbf{Arrays}: Inputs que aceptan arrays en formato JSON o separados por comas (por ejemplo, para métodos ICA)
\end{itemize}

\textbf{Validación y formateo de parámetros:}
\begin{itemize}
    \item \textbf{Validación de JSON}: Los campos tipo JSON permiten validación y formateo automático, mostrando mensajes de error o éxito en la interfaz.
    \item \textbf{Validación contextual}: Para plantillas asociadas a \textbf{board\_id} o \textbf{streamer\_type}, se verifica la presencia de campos requeridos y se sugieren mejoras.
    \item \textbf{Limpieza de metadatos}: Antes de enviar, los campos JSON eliminan claves auxiliares (como \textbf{\_COMMENTS} o campos que empiezan por \textbf{\_}).
    \item \textbf{Carga automática de plantillas}: Al cambiar el valor de campos selectores (como \textbf{board\_id}), se actualiza automáticamente la plantilla JSON asociada.
    \item \textbf{Botones de acción}: La interfaz incluye botones para cargar plantillas, formatear y validar JSON, con retroalimentación visual inmediata.
\end{itemize}

\textbf{Ejemplo de estructura de parámetros enviados:}
\begin{verbatim}
{
  "board_id": 1,
  "settings": {
    "gain": 24,
    "channels": [1,2,3,4]
  },
  "custom_template": {
    "param1": "value",
    "param2": 42
  }
}
\end{verbatim}
Los arrays pueden enviarse como \textbf{[1,2,3]} o \textbf{"1,2,3"} (se convierten automáticamente). Los campos JSON se validan y formatean antes de enviarse.

\subsection{Patrones arquitectónicos implementados}
\label{subsec:patrones_api_explorer}

\subsubsection{Patrón Singleton}
El \textbf{ApiExplorer} implementa un singleton thread-safe que garantiza:
\begin{itemize}
    \item \textbf{Estado global consistente}: Una única fuente de verdad para el registro de módulos
    \item \textbf{Compartición eficiente de recursos}: Reutilización de conexiones y pools de workers
    \item \textbf{Acceso thread-safe}: Sincronización mediante locks para operaciones críticas
\end{itemize}

\subsubsection{Patrón Worker Pool}
La gestión de WebSockets implementa un patrón worker pool que proporciona:
\begin{itemize}
    \item \textbf{Procesamiento distribuido}: Distribución de carga entre múltiples workers
    \item \textbf{Balanceo de carga automático}: Distribución round-robin de mensajes
    \item \textbf{Escalado adaptativo}: Configuración de workers basada en el tipo de módulo
\end{itemize}

\subsubsection{Patrón Observer}
El sistema de eventos implementa el patrón observer para:
\begin{itemize}
    \item \textbf{Difusión de eventos}: Distribución de eventos a múltiples clientes conectados
    \item \textbf{Sistema de callbacks}: Notificaciones automáticas de cambios de estado
    \item \textbf{Desacoplamiento}: Comunicación indirecta entre componentes
\end{itemize}

\subsubsection{Patrón Registry}
El registro dinámico de módulos utiliza el patrón registry para:
\begin{itemize}
    \item \textbf{Registro dinámico}: Adición de módulos y métodos en tiempo de ejecución
    \item \textbf{Configuración runtime}: Modificación de la estructura API sin reinicio
    \item \textbf{Extensibilidad}: Fácil integración de nuevos módulos
\end{itemize}

\subsection{Estructuras de datos}
\label{subsec:estructuras_datos_api_explorer}

\subsubsection{Información de módulos}

El sistema mantiene una estructura jerárquica para organizar los módulos registrados:

\begin{itemize}
    \item \textbf{Nivel módulo}: Metadatos del módulo (nombre, descripción, configuración)
    \item \textbf{Nivel grupo}: Agrupación lógica de métodos relacionados
    \item \textbf{Nivel método}: Especificación detallada de cada método API
    \item \textbf{Nivel parámetro}: Definición de tipos y validaciones de parámetros
\end{itemize}

\subsubsection{Métricas de rendimiento}

El sistema \textbf{WorkerStats} mantiene métricas detalladas por worker:
\begin{itemize}
    \item \textbf{Procesamiento de mensajes}: Conteo de mensajes procesados
    \item \textbf{Métricas de tiempo}: Tiempo promedio de procesamiento
    \item \textbf{Utilización de colas}: Porcentaje de utilización de colas de mensajes
    \item \textbf{Seguimiento de actividad}: Timestamp de última actividad
\end{itemize}

\subsection{Características de rendimiento}
\label{subsec:rendimiento_api_explorer}

\subsubsection{Manejo de datos de alta frecuencia}

Para datos de alta frecuencia como señales EEG (250-1000 Hz), el sistema implementa:
\begin{itemize}
    \item \textbf{Worker Pool Processing}: Procesamiento distribuido entre múltiples workers
    \item \textbf{Adaptive Throttling}: Throttling inteligente basado en frecuencia de datos
    \item \textbf{Batch Processing}: Agrupación de eventos para eficiencia de red
\end{itemize}

\subsubsection{Optimizaciones de UI}

Para mantener la responsividad de la interfaz:
\begin{itemize}
    \item \textbf{Batched DOM Updates}: Actualizaciones del DOM por lotes
    \item \textbf{Throttled Chart Updates}: Limitación de actualizaciones de gráficos
    \item \textbf{Incremental Table Updates}: Actualización incremental de tablas de datos
\end{itemize}

\subsubsection{Gestión de memoria}

El sistema implementa estrategias para evitar fugas de memoria:
\begin{itemize}
    \item \textbf{Automatic Cleanup}: Limpieza automática de recursos no utilizados
    \item \textbf{Connection Management}: Gestión del ciclo de vida de conexiones WebSocket
    \item \textbf{Buffer Limits}: Límites configurables para buffers de datos
\end{itemize}

\section{Diseño del módulo input\_collector}
\label{sec:diseno_input_collector}

\subsection{Propósito y responsabilidades}
\label{subsec:proposito_input_collector}

Input Collector es una aplicación integral diseñada para capturar, procesar, visualizar y registrar eventos de entrada en tiempo real desde múltiples dispositivos (teclado, ratón, controladores PS4). La aplicación implementa una \textbf{arquitectura modular basada en eventos} que soporta múltiples interfaces de usuario (GUI, CLI, Web API) y proporciona capacidades extensibles para diferentes tipos de colectores y modos de visualización.

\subsection{Arquitectura del módulo}
\label{subsec:arquitectura_input_collector}

La aplicación sigue una \textbf{arquitectura hexagonal (Ports and Adapters)} combinada con patrones \textbf{Mediator}, \textbf{Observer} y \textbf{Event-Driven Architecture}, organizando el código en capas bien definidas como se muestra en la figura~\ref{fig:arquitectura_input_collector}.

\begin{figure}[hp!]
\centering
\includegraphics[width=\textwidth,height=15cm,keepaspectratio]{imaxes/input_collector_arquitectura.pdf}
\caption{Arquitectura hexagonal del módulo input\_collector}
\label{fig:arquitectura_input_collector}
\end{figure}

La arquitectura se estructura en cinco capas principales:

\subsubsection{Interfaces Externas}
Múltiples puntos de acceso al sistema:
\begin{itemize}
    \item \textbf{GUI (PyQt5)}: Interfaz gráfica con tema oscuro y paneles divididos
    \item \textbf{CLI (Typer)}: Shell interactivo con monitoreo específico por colector
    \item \textbf{Web API (API Explorer)}: Control remoto con WebSocket en tiempo real
\end{itemize}

\subsubsection{Adaptadores}
Capa de traducción entre interfaces y servicios core:
\begin{itemize}
    \item \textbf{QtInteractionAdapter}: Adaptador thread-safe para operaciones UI
    \item \textbf{CLI Commands}: Comandos especializados para diferentes funcionalidades
    \item \textbf{API Integrations}: Integraciones web para datos y control
\end{itemize}

\subsubsection{Capa de Servicios}
Servicios singleton para lógica de negocio:
\begin{itemize}
    \item \textbf{InteractionService}: Control de aplicación expandido con thread-safety
    \item \textbf{DataAccessService}: Acceso a datos con categorización de callbacks
    \item \textbf{Core Event System}: Sistema de eventos framework-independiente
\end{itemize}

\subsubsection{Núcleo de Aplicación}
Coordinación central y gestión de componentes:
\begin{itemize}
    \item \textbf{App Coordinator}: Mediador central con integración al sistema de eventos
    \item \textbf{Collectors Manager}: Gestor con capabilities dinámicas
    \item \textbf{Processing}: Procesamiento y transformación de eventos
\end{itemize}

\subsubsection{Dominio}
Componentes fundamentales del sistema:
\begin{itemize}
    \item \textbf{Interfaces}: Contratos y estructuras de datos (InputData con 25+ campos)
    \item \textbf{Visualizers}: Componentes de visualización en tiempo real
    \item \textbf{Logging}: Sistema de registro a múltiples formatos
\end{itemize}

\subsection{Core Event System}
\label{subsec:core_event_system_input}

\subsubsection{Sistema de eventos centralizado}

El \textbf{Core Event System} actúa como backbone de comunicación, proporcionando un sistema de eventos framework-independiente:

\textbf{Características principales:}
\begin{itemize}
    \item \textbf{Arquitectura dirigida por eventos}: Comunicación desacoplada entre módulos
    \item \textbf{Thread-safe}: Múltiples locks granulares para diferentes operaciones
    \item \textbf{Aislamiento de errores}: Los fallos en callbacks no afectan otros suscriptores
    \item \textbf{Sincronización tardía}: Almacenamiento del último evento para suscriptores posteriores
    \item \textbf{Framework-independiente}: No depende de Qt u otros frameworks
\end{itemize}

\subsubsection{EventDispatcher singleton}

El \textbf{EventDispatcher} implementa un patrón singleton thread-safe que gestiona suscripciones y dispatch de eventos:

\textbf{Funcionalidades clave:}
\begin{itemize}
    \item \textbf{Suscripción dinámica}: Registro y desregistro de callbacks en tiempo de ejecución
    \item \textbf{Dispatch thread-safe}: Distribución segura de eventos entre hilos
    \item \textbf{Tipado de eventos}: Enumeración \textbf{EventType} para eventos del sistema
    \item \textbf{API simplificada}: Funciones utilitarias para uso común
\end{itemize}

\subsection{Colectores de entrada}
\label{subsec:colectores_entrada}

\subsubsection{CollectorManager con capabilities dinámicas}

El \textbf{CollectorManager} proporciona gestión centralizada de múltiples tipos de colectores:

\textbf{Capacidades soportadas:}
\begin{itemize}
    \item \textbf{Detección dinámica}: Capabilities específicas por colector
    \item \textbf{Switching en caliente}: Cambio de colector sin reinicio
    \item \textbf{Gestión de estado}: Tracking de estado activo y configuración
    \item \textbf{Interfaz uniforme}: API consistente independiente del dispositivo
\end{itemize}

\subsubsection{KeyboardMouseCollector}

Colector con múltiples modos de tracking para dispositivos tradicionales:

\textbf{Modos de tracking de ratón:}
\begin{itemize}
    \item \textbf{Delta}: Raw movement deltas para precisión máxima
    \item \textbf{Angle}: Dirección y ángulo (8-direccional) para análisis de patrones
    \item \textbf{Axis}: Direcciones cardinales (N/S/E/W) simplificadas
    \item \textbf{Joystick}: Dirección y velocidad (analog-like) para visualización
\end{itemize}

\subsubsection{PS4ControllerCollector con arquitectura event-driven}

El colector de controlador PS4 implementa una arquitectura completamente event-driven:

\textbf{Características arquitectónicas:}
\begin{itemize}
    \item \textbf{Event loop dedicado}: SDL2 event loop en hilo dedicado
    \item \textbf{Processing thread-safe}: Locks granulares y manejo de errores
    \item \textbf{Reconexión automática}: Manejo de desconexiones inalámbricas
    \item \textbf{Core event integration}: Dispatch automático al sistema de eventos
    \item \textbf{Detección de cambios}: Tracking de estados previos para triggers/sticks
\end{itemize}

\textbf{Capacidades del PS4 Controller:}
\begin{itemize}
    \item \textbf{Soporte analog completo}: 2 sticks analógicos y 2 triggers
    \item \textbf{14 botones}: Incluyendo D-pad, face buttons y system buttons
    \item \textbf{Conexión SDL2}: Soporte nativo via SDL2 GameController
    \item \textbf{Thread-safe events}: Procesamiento seguro entre hilos
\end{itemize}

\subsection{Sistema de visualización avanzada}
\label{subsec:visualizacion_input_collector}

\subsubsection{Visualizadores especializados}

El sistema incluye visualizadores específicos para cada tipo de dispositivo:

\textbf{KeyboardVisualizer:}
\begin{itemize}
    \item \textbf{Teclado virtual}: Representación visual con estados de teclas
    \item \textbf{Highlighting dinámico}: Resaltado de teclas activas
    \item \textbf{Layout configurable}: Soporte para diferentes distribuciones
\end{itemize}

\textbf{MouseVisualizer:}
\begin{itemize}
    \item \textbf{Visualización joystick}: Representación analógica del movimiento
    \item \textbf{Múltiples modos}: Adaptación visual según modo de tracking
    \item \textbf{Indicadores de botón}: Estado visual de botones del ratón
\end{itemize}

\subsubsection{PS4ControllerVisualizer con interfaz de paneles divididos}

El visualizador del controlador PS4 implementa una interfaz avanzada de paneles divididos:

\textbf{Estructura de paneles:}
\begin{itemize}
    \item \textbf{Panel superior}: Visualización del controlador con image overlay
    \item \textbf{Panel medio}: Estado en tiempo real con valores analógicos
    \item \textbf{Panel inferior}: Log de eventos con timestamps y colores
\end{itemize}

\textbf{Características avanzadas:}
\begin{itemize}
    \item \textbf{Updates thread-safe}: Señales Qt para actualizaciones desde hilos
    \item \textbf{Edit mode avanzado}: Sistema snapshot/commit para posiciones
    \item \textbf{Múltiples modos}: Image overlay + dibujado directo
    \item \textbf{Logging en tiempo real}: Panel de logs con colores y timestamps
\end{itemize}

\subsection{Servicios de datos y control}
\label{subsec:servicios_input_collector}

\subsubsection{DataAccessService con categorización}

El \textbf{DataAccessService} proporciona acceso estructurado a datos con categorización avanzada:

\textbf{Sistema de callbacks categorizados:}
\begin{itemize}
    \item \textbf{API callbacks}: Para consumidores externos con throttling
    \item \textbf{Local callbacks}: Para visualizadores locales sin throttling
    \item \textbf{Categorización por evento}: Separación por tipo (key, mouse\_movement, controller)
    \item \textbf{Gestión de memoria}: Límites de historial configurables
\end{itemize}

\textbf{Thread-safety mejorado:}
\begin{itemize}
    \item \textbf{Locks granulares}: Diferentes locks para operaciones específicas
    \item \textbf{Estado metadata}: Tracking de información del colector activo
    \item \textbf{Patrones reactivos}: Sistema de callbacks para consumers externos
\end{itemize}

\subsubsection{InteractionService expandido}

El \textbf{InteractionService} gestiona el control de la aplicación con capacidades expandidas:

\textbf{Operaciones soportadas:}
\begin{itemize}
    \item \textbf{Control de grabación}: Start/stop con sincronización de estado
    \item \textbf{Switching de colectores}: Cambio dinámico con preservación de estado
    \item \textbf{Configuración de tracking}: Múltiples modos y parámetros
    \item \textbf{Gestión de UI}: Edit mode y configuración visual
\end{itemize}

\subsubsection{QtInteractionAdapter para thread-safety}

El adaptador Qt garantiza operaciones thread-safe entre API y GUI:

\textbf{Señales Qt implementadas (7 señales):}
\begin{itemize}
    \item \textbf{recording\_state\_changed}: Estado de grabación
    \item \textbf{collector\_switched}: Cambio de colector activo
    \item \textbf{mouse\_mode\_changed}: Modo de tracking del ratón
    \item \textbf{edit\_mode\_changed}: Estado del modo de edición
    \item \textbf{tracked\_keys\_changed}: Configuración de teclas tracked
    \item \textbf{mouse\_buttons\_changed}: Configuración de botones del ratón
    \item \textbf{ps4\_visualization\_mode\_changed}: Modo de visualización PS4
\end{itemize}

\subsection{Interfaz de línea de comandos}
\label{subsec:cli_input_collector}

\subsubsection{CLI avanzado con monitoreo específico}

El sistema CLI implementa capacidades avanzadas de monitoreo:

\textbf{Componentes principales:}
\begin{itemize}
    \item \textbf{CLIApp}: Aplicación principal con Typer y auto-completion
    \item \textbf{CLIShell}: Shell interactivo estilo REPL
    \item \textbf{PS4ControllerMonitor}: Monitor dedicado para controladores PS4
    \item \textbf{Commands especializados}: Módulos para diferentes funcionalidades
\end{itemize}

\subsubsection{Monitor de controlador PS4}

El monitor especializado para PS4 proporciona visualización en tiempo real:

\textbf{Características del monitor:}
\begin{itemize}
    \item \textbf{Rich console interface}: Tablas y paneles en tiempo real
    \item \textbf{Cross-platform key detection}: ESC/Q para salir en Windows/Unix
    \item \textbf{Event-driven architecture}: Uso del core event system
    \item \textbf{Smart collector detection}: Monitoreo específico según colector activo
\end{itemize}

\textbf{Estructura de comandos:}
\begin{itemize}
    \item \textbf{Data access}: \textbf{input-collector data active-keys}
    \item \textbf{Control}: \textbf{input-collector interaction recording start}
    \item \textbf{Monitoreo}: \textbf{input-collector monitor controller}
    \item \textbf{Configuración}: \textbf{input-collector interaction mouse set-mode joystick}
\end{itemize}

\subsection{Sistema de logging y procesamiento}
\label{subsec:logging_procesamiento_input}

\subsubsection{Logging multi-formato}

El \textbf{LoggingManager} coordina registro simultáneo a múltiples formatos:

\textbf{Formatos soportados:}
\begin{itemize}
    \item \textbf{JSON Logger}: Formato JSON estructurado para análisis programático
    \item \textbf{CSV Logger}: Formato tabular con headers dinámicos
    \item \textbf{Headers dinámicos CSV}: Generación automática basada en tracked keys
    \item \textbf{Thread-safe signals}: Comunicación mediante señales Qt
\end{itemize}

\subsubsection{EventProcessor extendido}

El procesador de eventos maneja formateo específico por colector:

\textbf{Características de procesamiento:}
\begin{itemize}
    \item \textbf{Formateo específico}: Diferente handling para keyboard/mouse vs PS4
    \item \textbf{Detección de eventos especiales}: F9 para toggle recording, eventos de sistema
    \item \textbf{Procesamiento de ejes}: Handling especializado para joysticks y triggers
    \item \textbf{Thread-safe}: Procesamiento seguro con señales Qt
\end{itemize}

\subsection{Integración con APIs web}
\label{subsec:integracion_api_input}

\subsubsection{API Integration con core events}

El sistema de integración API proporciona exposición web con eventos core:

\textbf{Componentes de integración:}
\begin{itemize}
    \item \textbf{InputCollectorAPI}: API estática para control de aplicación
    \item \textbf{InputDataAPI}: API estática para acceso a datos
    \item \textbf{Core event integration}: WebSocket broadcasts automáticos
    \item \textbf{PS4 state tracking}: Estado especializado para controladores
\end{itemize}

\textbf{Características web:}
\begin{itemize}
    \item \textbf{Enhanced fallbacks}: Operación sin web server mejorada
    \item \textbf{Thread-safe operations}: Todas las operaciones API son thread-safe
    \item \textbf{Real-time broadcasting}: Eventos automáticos via WebSocket
    \item \textbf{State synchronization}: Sincronización bidireccional GUI-API
\end{itemize}

\subsection{Patrones arquitectónicos implementados}
\label{subsec:patrones_input_collector}

\subsubsection{Event-Driven Architecture}
El sistema completo se basa en comunicación dirigida por eventos:
\begin{itemize}
    \item \textbf{Core Event System}: Backbone de comunicación framework-independiente
    \item \textbf{Pub/Sub pattern}: Desacoplamiento entre componentes
    \item \textbf{Automatic broadcasting}: Eventos automáticos para APIs y CLI
    \item \textbf{Late subscriber synchronization}: Estabilidad para suscriptores tardíos
\end{itemize}

\subsubsection{Patrón Mediator}
El \textbf{AppCoordinator} actúa como mediador central:
\begin{itemize}
    \item \textbf{Central coordination}: Gestión de estado centralizada
    \item \textbf{Core event integration}: Comunicación con sistemas externos
    \item \textbf{Thread-safe coordination}: Uso de señales Qt para operaciones UI
    \item \textbf{Error handling}: Manejo de errores en coordinación
\end{itemize}

\subsubsection{Patrón Singleton}
Servicios críticos implementan singleton thread-safe:
\begin{itemize}
    \item \textbf{Thread-safe implementation}: Double-checked locking
    \item \textbf{Granular locking}: Locks específicos para operaciones
    \item \textbf{Error isolation}: Aislamiento de fallos para estabilidad
    \item \textbf{Resource management}: Gestión centralizada de recursos
\end{itemize}

\subsubsection{Patrón Strategy}
Adaptación dinámica según contexto:
\begin{itemize}
    \item \textbf{Collector capabilities}: Capacidades dinámicas por dispositivo
    \item \textbf{Adaptive UI}: Interfaz adaptativa según colector activo
    \item \textbf{Multiple tracking modes}: Modos intercambiables de tracking
    \item \textbf{Platform-specific handling}: Adaptación por sistema operativo
\end{itemize}

\subsection{Características de rendimiento}
\label{subsec:rendimiento_input_collector}

\subsubsection{Thread safety avanzado}

El sistema implementa thread safety a múltiples niveles:
\begin{itemize}
    \item \textbf{Multiple granular locks}: Locks específicos en servicios y collectors
    \item \textbf{Qt signal/slot threading}: Updates UI thread-safe
    \item \textbf{SDL2 event loop threading}: Procesamiento PS4 en hilo dedicado
    \item \textbf{Error isolation}: Aislamiento de errores en callbacks
\end{itemize}

\subsubsection{Optimizaciones de entrada}

Para dispositivos de alta frecuencia como controladores:
\begin{itemize}
    \item \textbf{Event-driven processing}: Procesamiento basado en eventos reales
    \item \textbf{Change detection}: Tracking de cambios para reducir overhead
    \item \textbf{Efficient polling}: 1ms polling para SDL2 con optimizaciones
    \item \textbf{Memory management}: Gestión eficiente de buffers de eventos
\end{itemize}

\subsubsection{Configurabilidad y extensibilidad}

El sistema proporciona alta configurabilidad:
\begin{itemize}
    \item \textbf{Dynamic capabilities}: Adaptación automática según dispositivo
    \item \textbf{Runtime configuration}: Cambios sin reinicio de aplicación
    \item \textbf{Modular architecture}: Fácil adición de nuevos colectores
    \item \textbf{Plugin-ready design}: Preparado para extensiones futuras
\end{itemize}

\section{Diseño del módulo eeg\_processor}
\label{sec:diseno_eeg_processor}

\subsection{Propósito y responsabilidades}
\label{subsec:proposito_eeg_processor}

EEG Processor es un framework modular para procesamiento de señales EEG en tiempo real que soporta múltiples dispositivos de hardware (OpenBCI, g.tec, Muse, ANT Neuro, Neurosity), análisis avanzado (ICA, wavelet transforms, MNE-Python), y múltiples interfaces (GUI Qt5, CLI, REST API). Implementa una \textbf{arquitectura orientada a servicios} con \textbf{patrones dirigidos por eventos} para procesamiento de neuroseñales.

\subsection{Arquitectura del módulo}
\label{subsec:arquitectura_eeg_processor}

La aplicación sigue una \textbf{arquitectura orientada a servicios (SOA)} combinada con patrones \textbf{dirigidos por eventos}, \textbf{Singleton} y \textbf{MVC}, organizando el código en capas especializadas como se muestra en la figura~\ref{fig:arquitectura_eeg_processor}.

\begin{figure}[hp!]
\centering
\includegraphics[width=\textwidth,height=15cm,keepaspectratio]{imaxes/arquitectura_eeg_processor.pdf}
\caption{Arquitectura SOA del módulo eeg\_processor}
\label{fig:arquitectura_eeg_processor}
\end{figure}

La arquitectura se estructura en cinco capas principales:

\subsubsection{Interfaces Externas}
Múltiples puntos de acceso al sistema:
\begin{itemize}
    \item \textbf{GUI (Qt5 MVC)}: Interfaz basada en paneles con controladores especializados
    \item \textbf{CLI (Typer)}: Shell interactivo con salida Rich y monitoreo tiempo real
    \item \textbf{REST API (WebSocket+HTTP)}: Control remoto con streaming de alta frecuencia
\end{itemize}

\subsubsection{Adaptadores}
Puentes thread-safe entre frameworks:
\begin{itemize}
    \item \textbf{QtInteractionAdapter}: Adaptador thread-safe con ThreadPoolExecutor
    \item \textbf{CLI Commands}: Comandos especializados para datos, monitoreo e interacción
    \item \textbf{API Integrations}: Integración WebSocket con modo dual (directo/multiproceso)
\end{itemize}

\subsubsection{Capa de Servicios}
Servicios singleton para lógica de negocio especializada:
\begin{itemize}
    \item \textbf{CollectorService}: Gestión singleton de colección EEG con soporte multi-preset
    \item \textbf{Data Access Service}: Acceso centralizado con gestión de buffer thread-safe
    \item \textbf{Advanced Analysis Services}: ICA, MNE, wavelet con multiprocesamiento
\end{itemize}

\subsubsection{Controladores UI}
Arquitectura MVC basada en paneles:
\begin{itemize}
    \item \textbf{Panel-Based MVC}: Controladores coordinando múltiples widgets especializados
    \item \textbf{MNE Panel Controller}: Conversión y visualización MNE-Python asíncrona
    \item \textbf{Advanced Panel Controller}: ICA, wavelet, filtrado avanzado
\end{itemize}

\subsubsection{Núcleo}
Componentes fundamentales del framework:
\begin{itemize}
    \item \textbf{Event System}: Sistema de eventos con 40+ tipos y tracking histórico
    \item \textbf{Signal Processing}: Motor con múltiples algoritmos y monitoreo de estabilidad
    \item \textbf{Hardware Abstraction}: Capa BrainFlow para 20+ fabricantes
\end{itemize}

\subsection{Core Event System}
\label{subsec:core_event_system_eeg}

\subsubsection{Sistema de eventos especializado}

El sistema de eventos del módulo EEG extiende el concepto básico con características específicas para procesamiento de neuroseñales:

\textbf{Características especializadas:}
\begin{itemize}
    \item \textbf{40+ tipos de eventos}: CONNECTION, PROCESSING, VISUALIZATION, MNE, ICA, WAVELET
    \item \textbf{Seguimiento histórico}: 1000 eventos con rotación automática para análisis
    \item \textbf{Operaciones thread-safe}: RLock protection para operaciones concurrentes
    \item \textbf{Suscriptores tardíos}: Entrega automática del último evento por tipo
    \item \textbf{Aislamiento de errores}: Fallos en callbacks no comprometen otros suscriptores
\end{itemize}

\subsubsection{Event payload estructurado}

Los eventos incluyen dataclass \textbf{Event} con timestamp, source y data payload para trazabilidad completa y debugging avanzado.

\subsection{Capa de abstracción de hardware}
\label{subsec:abstraccion_hardware_eeg}

\subsubsection{EEGCollector con soporte multi-dispositivo}

El \textbf{EEGCollector} proporciona una interfaz unificada para 20+ fabricantes de dispositivos EEG:

\textbf{Ejemplos de dispositivos soportados:}
\begin{itemize}
    \item \textbf{OpenBCI}: Cyton Board, Ganglion Board con configuración personalizable
    \item \textbf{g.tec}: Unicorn Board con mapeo automático de canales
    \item \textbf{Muse}: Muse 2, Muse S con detección de tipo de señal
    \item \textbf{ANT Neuro}: eego mylab con soporte profesional
    \item \textbf{Neurosity}: Notion 1, Notion 2 con métricas cognitivas
    \item \textbf{BrainBit}: BrainBit, BrainBit BLED con conectividad Bluetooth
\end{itemize}

\subsubsection{Configuración avanzada de dispositivos}

\textbf{BrainFlowBoardConfiguration} gestiona configuración específica por fabricante:
\begin{itemize}
    \item \textbf{Soporte multi-preset}: Streams DEFAULT, AUXILIARY, ANCILLARY
    \item \textbf{Inserción de marcadores tiempo real}: Sincronización de microsegundo
    \item \textbf{Detección automática de canales}: EEG, EMG, ECG, acelerómetro, giroscopio
    \item \textbf{Placas de streaming}: Configuración de placa maestra con validación de preset
    \item \textbf{Operaciones thread-safe}: Gestión de buffer y recuperación de datos
\end{itemize}

\subsection{Motor de procesamiento de señales}
\label{subsec:motor_procesamiento_eeg}

\subsubsection{SignalProcessor con algoritmos avanzados}

El motor de procesamiento implementa múltiples familias de algoritmos con optimizaciones específicas:

\textbf{Filtrado avanzado:}
\begin{itemize}
    \item \textbf{Múltiples implementaciones}: Butterworth, Chebyshev Tipo 1, Bessel
    \item \textbf{Procesamiento fase cero}: Eliminación de distorsión de fase
    \item \textbf{Monitoreo de estabilidad}: Ajuste automático para datos sintéticos
    \item \textbf{Filtros especializados}: Notch (50/60 Hz), detrend, pasa-banda configurable
\end{itemize}

\textbf{Análisis avanzado:}
\begin{itemize}
    \item \textbf{Análisis wavelet}: 15+ familias (Daubechies, Biorthogonal, Coiflets, Symlets)
    \item \textbf{FastICA}: Configuración de 1-16 componentes con selección de canales
    \item \textbf{Métricas EEG}: Potencias de banda (Delta, Theta, Alpha, Beta, Gamma)
    \item \textbf{Métricas cognitivas}: Puntuaciones de atención, concentración, relajación
    \item \textbf{Evaluación de calidad}: Sistema de 5 niveles con detección de saturación
\end{itemize}

\subsection{Servicios especializados}
\label{subsec:servicios_especializados_eeg}

\subsubsection{Servicios ICA con multiprocesamiento}

Los servicios ICA implementan procesamiento distribuido para análisis de componentes independientes:

\textbf{ICAProcessService:}
\begin{itemize}
    \item \textbf{Procesamiento multiproceso}: Integración con señales Qt para UI responsiva
    \item \textbf{Configuración flexible}: 1-16 componentes con selección de canales
    \item \textbf{Monitoreo tiempo real}: ICARealtimeService con ventana deslizante
    \item \textbf{Exportación por lotes}: ExportService con seguimiento de progreso
\end{itemize}

\subsubsection{Servicios MNE-Python}

Integración completa con MNE-Python para análisis neurocientífico avanzado:

\textbf{MNEConverterService:}
\begin{itemize}
    \item \textbf{Conversión asíncrona}: BrainFlow → MNE con múltiples modos
    \item \textbf{Modo flexible}: Conversión de todos los canales para análisis completo
    \item \textbf{Modo rendimiento}: Canales seleccionados para optimización
    \item \textbf{Modo tiempo real}: Ventana deslizante para procesamiento continuo
\end{itemize}

\textbf{MNEVisualizerService:}
\begin{itemize}
    \item \textbf{Visualización avanzada}: PSD, topomaps, espectrogramas, datos en bruto
    \item \textbf{Calidad publicación}: Gráficos de alta resolución (200 DPI)
    \item \textbf{Procesamiento tiempo real}: Sesiones configurables con updates dinámicos
    \item \textbf{MNEProcessManager}: Coordinación multiproceso para operaciones CPU-intensivas
\end{itemize}

\subsection{Arquitectura UI basada en paneles}
\label{subsec:arquitectura_ui_eeg}

\subsubsection{Sistema MVC especializado}

La interfaz implementa una arquitectura MVC basada en paneles con controladores especializados:

\textbf{Paneles principales:}
\begin{itemize}
    \item \textbf{ConnectionPanel}: Selección de placa con validación dinámica
    \item \textbf{ProcessingPanel}: Filtros tiempo real con visualización de mapa cerebral
    \item \textbf{MNE Panel}: Conversión multi-modo con visualización avanzada
    \item \textbf{Advanced Panel}: ICA, wavelet, filtrado avanzado
\end{itemize}

\subsubsection{ConnectionPanel con validación dinámica}

\textbf{Características del panel de conexión:}
\begin{itemize}
    \item \textbf{Selección inteligente}: Filtrado y categorización por fabricante
    \item \textbf{Validación dinámica}: Retroalimentación visual para campos requeridos/opcionales
    \item \textbf{Configuración streaming}: Validación de preset con prevención de errores
    \item \textbf{Información tiempo real}: Mapeo de canales, frecuencia de muestreo, detalles de hardware
\end{itemize}

\subsubsection{ProcessingPanel con filtros tiempo real}

\textbf{Capacidades de procesamiento:}
\begin{itemize}
    \item \textbf{Aplicación tiempo real}: Filtros durante procesamiento activo sin reinicio
    \item \textbf{Filtrado integral}: Pasabanda (1-200 Hz), Notch (50/60 Hz), Detrend
    \item \textbf{Métricas EEG}: Cálculo de potencias de banda con análisis estadístico
    \item \textbf{Mapa cerebral}: Visualización optimizada con datos de streaming
\end{itemize}

\subsubsection{MNE Panel con principios SOLID}

El panel MNE implementa principios SOLID con composición de widgets:

\textbf{MNEPanelController:}
\begin{itemize}
    \item \textbf{Conversión multi-modo}: Flexible, rendimiento, tiempo real
    \item \textbf{Visualización avanzada}: PSD, topomaps, espectrogramas coordinados
    \item \textbf{Sesiones tiempo real}: Actualizaciones de parámetros sin reinicio
    \item \textbf{Integración filtros}: Coordinación fluida con panel de procesamiento
\end{itemize}

\textbf{Widgets especializados:}
\begin{itemize}
    \item \textbf{MNEConversionWidget}: Interfaz multi-modo con advertencias de rendimiento
    \item \textbf{MNEVisualizationWidget}: Controles con actualizaciones tiempo real
    \item \textbf{ChannelSelectionDialog}: Selección basada en cuadrícula con optimización
\end{itemize}

\subsection{Sistema de visualización científica}
\label{subsec:visualizacion_cientifica_eeg}

\subsubsection{MultiSignalVisualizer~\ref{fig:multi_signal_viz}}

Visualización biosignal multi-página con detección inteligente:

\textbf{Características avanzadas:}
\begin{itemize}
    \item \textbf{Detección inteligente}: EEG, EMG, ECG, EOG, EDA, PPG, acelerómetro, giroscopio
    \item \textbf{Arquitectura carga perezosa}: Caché para optimización de rendimiento
    \item \textbf{Interfaz multi-pestaña}: Separación por modalidad de señal
    \item \textbf{Integración tiempo real}: Coordinación con DataAccessService
    \item \textbf{Optimización rendimiento}: Actualizaciones selectivas, operaciones eficientes en memoria
\end{itemize}

\begin{figure}[hp!]
\centering
\includegraphics[width=\textwidth,height=7cm,keepaspectratio]{imaxes/multi_signal_viz.png}
\caption{MultiSignalVisualizer}
\label{fig:multi_signal_viz}
\end{figure}

\subsubsection{BrainMapVisualizer}

Mapeo topográfico cerebral con estándares internacionales:

\textbf{Características del mapa cerebral:}
\begin{itemize}
    \item \textbf{Posicionamiento 10-20/10-10}: Posiciones estándar de electrodos
    \item \textbf{Calidad tiempo real}: Indicadores codificados por color (5 niveles)
    \item \textbf{Regiones anatómicas}: Superposiciones frontal, central, parietal, occipital
    \item \textbf{Compatibilidad G.tec}: Mapeo automático de canales para Unicorn
    \item \textbf{Elementos interactivos}: Etiquetas de electrodos con leyenda de calidad
\end{itemize}

En la figura ~\ref{fig:brain_map_viz} se puede observar la disposicion de los electrodos en el brain map junto a la calidad de la señal.
En la figura ~\ref{fig:g_tec_suite} se puede observar la misma información en la propia suite de g.tec

\begin{figure}[hp!]
\centering
\includegraphics[width=\textwidth,height=7cm,keepaspectratio]{imaxes/brain_map_viz.png}
\caption{Brain Map Visualizer}
\label{fig:brain_map_viz}
\end{figure}
\begin{figure}[hp!]
\centering
\includegraphics[width=\textwidth,height=5cm,keepaspectratio]{imaxes/g_tec_suite.png}
\caption{G.Tec Suite Unicorn Hybrid Black}
\label{fig:g_tec_suite}
\end{figure}

\subsection{CLI avanzado con Rich console}
\label{subsec:cli_avanzado_eeg}

\subsubsection{Sistema CLI modular}

El CLI implementa capacidades avanzadas con framework Typer:

\textbf{Componentes principales:}
\begin{itemize}
    \item \textbf{CLIApp}: Aplicación Typer con autocompletado y registro de comandos
    \item \textbf{CLIShell}: Shell interactivo estilo REPL con historial
    \item \textbf{Módulos especializados}: DataCommands, MonitorCommands, InteractionCommands
\end{itemize}

\subsubsection{Comandos especializados}

\textbf{Estructura de comandos:}
\begin{itemize}
    \item \textbf{Data access}: \textbf{data current-data --max-samples 1000}
    \item \textbf{Control}: \textbf{control connect-board --board-id 1 --preset 0}
    \item \textbf{Monitoreo}: \textbf{monitor realtime --channels 1,2,3,4}
    \item \textbf{Procesamiento}: \textbf{control apply-filters --bandpass-low 1 --bandpass-high 50}
\end{itemize}

\textbf{Características avanzadas:}
\begin{itemize}
    \item \textbf{Salida Rich}: Tablas, paneles, visualizaciones tiempo real
    \item \textbf{Soporte multiplataforma}: Compatibilidad Windows/Linux/macOS
    \item \textbf{Shell interactivo}: Interfaz REPL con historial de comandos
\end{itemize}

\subsection{Integración con APIs web}
\label{subsec:integracion_api_eeg}

\subsubsection{API con WebSocket}

El sistema API proporciona streaming de alta frecuencia con modo dual:

\textbf{Componentes de integración:}
\begin{itemize}
    \item \textbf{EEGDataAPI}: API estática para acceso a datos y metadatos
    \item \textbf{EEGProcessorInteractionAPI}: API estática para control del sistema
    \item \textbf{Dos modos de Streaming}: Directo (baja latencia ~15ms), Multiproceso (más latencia, pero para procesamientos demandantes)
    \item \textbf{WebSocket tiempo real}: Difusión de payload estructurado para datos EEG
\end{itemize}

\textbf{Capacidades API:}
\begin{itemize}
    \item \textbf{Gestión de placas}: Conexión, configuración, detección de dispositivos
    \item \textbf{Control de procesamiento}: Filtros, métricas, análisis avanzado tiempo real
    \item \textbf{Grabación}: Inicio/parada con inserción de marcadores
    \item \textbf{Análisis}: ICA, wavelet, conversión MNE, exportación de resultados
\end{itemize}

\subsection{Patrones arquitectónicos implementados}
\label{subsec:patrones_eeg_processor}

\subsubsection{Arquitectura Orientada a Servicios (SOA)}
El sistema se organiza en servicios especializados con responsabilidades claras:
\begin{itemize}
    \item \textbf{Servicios singleton}: Inicialización thread-safe con verificación doble
    \item \textbf{Coordinación de servicios}: InteractionService como orquestador central
    \item \textbf{Límites claros}: Interfaces bien definidas para cada área funcional
\end{itemize}

\subsubsection{Modelo-Vista-Controlador (MVC)}
La UI implementa MVC especializado para aplicaciones científicas:
\begin{itemize}
    \item \textbf{MVC basado en paneles}: Controladores coordinando múltiples widgets
    \item \textbf{Controladores especializados}: MNEPanelController, AdvancedPanelController
    \item \textbf{Comunicación señal/slot}: Arquitectura Qt para UI responsiva
    \item \textbf{Gestión de estado}: Consistencia a través de contextos operacionales
\end{itemize}

\subsubsection{Multiprocesamiento y Multihilo}
Arquitectura concurrente para operaciones CPU-intensivas:
\begin{itemize}
    \item \textbf{Operaciones CPU-intensivas}: ICA, MNE en procesos separados
    \item \textbf{Hilos en segundo plano}: Colección de datos, procesamiento tiempo real
    \item \textbf{Comunicación thread-safe}: Señales Qt, colas multiproceso
    \item \textbf{Gestión de ciclo de vida}: Limpieza apropiada y optimización de memoria
\end{itemize}

\subsubsection{Patrones Factory y Strategy}
Adaptación dinámica según contexto y dispositivo:
\begin{itemize}
    \item \textbf{Factory configuración}: Creación dinámica de parámetros por dispositivo
    \item \textbf{Strategy procesamiento}: Múltiples tipos de filtros y algoritmos
    \item \textbf{Strategy visualización}: UI adaptativa basada en capacidades del dispositivo
    \item \textbf{Strategy detección}: Reconocimiento inteligente de tipos de canal
\end{itemize}

\subsection{Características de rendimiento}
\label{subsec:rendimiento_eeg_processor}

\subsubsection{Optimizaciones tiempo real}

Para procesamiento EEG de alta frecuencia (250-1000 Hz):
\begin{itemize}
    \item \textbf{Pipelines optimizados}: Renderizado de 50-100ms para actualizaciones fluidas
    \item \textbf{Gestión adaptiva de timeout}: Basada en complejidad de procesamiento
    \item \textbf{Operaciones eficientes en memoria}: Gestión automática de buffer
    \item \textbf{Monitoreo de rendimiento}: Optimización de uso de CPU con advertencias
\end{itemize}

\subsubsection{Escalabilidad y concurrencia}

\begin{itemize}
    \item \textbf{Soporte multiproceso}: Operaciones ICA y MNE en procesos separados
    \item \textbf{Thread safety avanzado}: Locks granulares y comunicación por señales
    \item \textbf{Gestión de recursos}: Pools de hilos y limpieza automática
    \item \textbf{Streaming adaptativo}: Modo dual para diferentes requisitos de latencia
\end{itemize}

\subsubsection{Características educativas}

El sistema incluye características específicas para investigación y educación:
\begin{itemize}
    \item \textbf{Sistema de ayuda integral}: 25+ temas con antecedentes teóricos
    \item \textbf{Orientación contextual}: Recomendaciones de parámetros basadas en contexto
    \item \textbf{Complejidad progresiva}: Configuración básica → investigación avanzada
    \item \textbf{Documentación integrada}: Ayuda con mejores prácticas
\end{itemize}

\section{Integración entre módulos (placeholder, no definitivo)}
\label{sec:integracion_modulos}

\subsection{Comunicación inter-módulo}
\label{subsec:comunicacion_inter_modulo}

\subsubsection{Lab Streaming Layer (LSL)}

Para sincronización temporal de alta precisión entre módulos, el framework utiliza Lab Streaming Layer como protocolo estándar:

\textbf{Características de LSL:}
\begin{itemize}
    \item \textbf{Sincronización temporal}: Precisión de microsegundos para correlación de eventos
    \item \textbf{Streams especializados}: Input events (irregular), EEG data (250-1000 Hz)
    \item \textbf{Descubrimiento automático}: Auto-detección de streams en la red local
    \item \textbf{Buffering inteligente}: Gestión automática de buffers para diferentes frecuencias
    \item \textbf{Recuperación de timestamp}: Corrección de deriva de reloj entre dispositivos
\end{itemize}

\subsubsection{API REST común via api\_explorer}

Los módulos se registran dinámicamente en api\_explorer para coordinación centralizada:

\textbf{Registro modular:}
\begin{itemize}
    \item \textbf{input\_collector}: APIs para captura de eventos y control de dispositivos
    \item \textbf{eeg\_processor}: APIs para procesamiento de señales y análisis avanzado
    \item \textbf{Coordinación cruzada}: Inicio/parada sincronizado de sesiones de grabación
    \item \textbf{Estado compartido}: Información de estado entre módulos vía WebSocket
\end{itemize}

\subsection{Flujo de datos sincronizado}
\label{subsec:flujo_datos_sincronizado}

\textbf{Secuencia de sincronización:}
\begin{enumerate}
    \item \textbf{Inicialización}: Establecimiento de streams LSL en ambos módulos
    \item \textbf{Calibración temporal}: Sincronización de relojes y corrección de deriva
    \item \textbf{Inicio coordinado}: Comando simultáneo de grabación via api\_explorer
    \item \textbf{Streaming continuo}: Datos EEG y eventos de entrada con timestamps LSL
    \item \textbf{Correlación temporal}: Asociación de eventos con señales EEG
    \item \textbf{Finalización sincronizada}: Parada coordinada y exportación de datasets
\end{enumerate}

\section{Consideraciones de implementación}
\label{sec:consideraciones_implementacion}

\subsection{Thread safety y concurrencia}
\label{subsec:thread_safety}

\subsubsection{Estrategias de sincronización}

El framework implementa múltiples estrategias de sincronización adaptadas a diferentes contextos:

\begin{itemize}
    \item \textbf{Granular locking}: Locks específicos para operaciones críticas diferentes
    \item \textbf{Qt Signals/Slots}: Comunicación thread-safe entre hilos de worker y UI
    \item \textbf{Immutable data structures}: Uso de estructuras inmutables cuando es posible
    \item \textbf{Actor model}: En WebSocket workers para aislamiento de procesamiento
    \item \textbf{Event-driven communication}: Sistema de eventos framework-independiente
\end{itemize}

\subsubsection{Gestión de recursos thread-safe}

Los recursos críticos se gestionan mediante patrones thread-safe específicos:
\begin{itemize}
    \item \textbf{Resource pools}: Para buffers de datos de alta frecuencia
    \item \textbf{Connection management}: Gestión de conexiones de dispositivos
    \item \textbf{Memory management}: Limpieza automática de recursos no utilizados
    \item \textbf{Timeout handling}: Gestión de timeouts para operaciones bloqueantes
\end{itemize}

\subsection{Optimizaciones de rendimiento}
\label{subsec:optimizaciones_rendimiento}

\subsubsection{Gestión eficiente de memoria}

Para el manejo de datos de alta frecuencia, se implementan optimizaciones específicas:

\begin{itemize}
    \item \textbf{Circular buffers}: Para datos continuos con overwrite automático
    \item \textbf{Memory pooling}: Pre-asignación de chunks para evitar allocations frecuentes
    \item \textbf{Copy-on-write}: Optimización para compartir datos entre componentes
    \item \textbf{Garbage collection optimization}: Estrategias para minimizar GC pressure
\end{itemize}

\subsubsection{Optimizaciones de procesamiento}

\begin{itemize}
    \item \textbf{Vectorización NumPy}: Uso intensivo de operaciones vectorizadas
    \item \textbf{Paralelización}: Multiprocesamiento para operaciones CPU-intensivas
    \item \textbf{Caching inteligente}: Caché de resultados de computaciones costosas
    \item \textbf{Lazy evaluation}: Evaluación perezosa para operaciones opcionales
\end{itemize}

\subsection{Manejo de errores y recuperación}
\label{subsec:manejo_errores}

\subsubsection{Estrategias de recuperación}

El sistema implementa múltiples niveles de manejo de errores:

\begin{itemize}
    \item \textbf{Circuit breaker pattern}: Para protección contra fallos repetidos
    \item \textbf{Graceful degradation}: Funcionamiento reducido ante fallos parciales
    \item \textbf{Automatic retry}: Reintentos automáticos con backoff exponencial
    \item \textbf{Error isolation}: Aislamiento de errores para evitar propagación
    \item \textbf{Health monitoring}: Monitoreo continuo del estado del sistema
\end{itemize}

\subsubsection{Logging y diagnóstico}

\begin{itemize}
    \item \textbf{Structured logging}: Logs estructurados para análisis automatizado
    \item \textbf{Performance metrics}: Métricas de rendimiento para optimización
    \item \textbf{Error tracking}: Seguimiento detallado de errores y contexto
    \item \textbf{Debug modes}: Modos de debugging para desarrollo y producción
\end{itemize}

\section{Resumen del capítulo}
\label{sec:resumen_analisis_diseno}

Este capítulo ha presentado el análisis y diseño completo del framework desarrollado, abarcando:

\begin{enumerate}
    \item \textbf{Análisis de requisitos exhaustivo}: 24 requisitos funcionales y 17 no funcionales que guían el desarrollo del sistema, estableciendo las bases para un framework robusto y escalable.

    \item \textbf{Arquitectura general orientada a servicios}: El framework implementa una arquitectura SOA con comunicación dirigida por eventos, proporcionando:
    \begin{itemize}
        \item Separación clara de responsabilidades entre módulos
        \item Comunicación desacoplada mediante múltiples protocolos (Core Events, REST, WebSocket, LSL)
        \item Soporte multi-interfaz simultáneo (GUI, CLI, API REST)
        \item Extensibilidad y mantenibilidad a largo plazo
    \end{itemize}

    \item \textbf{Diseño detallado de tres (placeholder) módulos principales}:
    \begin{itemize}
        \item \textbf{api\_explorer}: Framework web modular con arquitectura singleton, worker pools para alta concurrencia, y frontend JavaScript basado en managers especializados
        \item \textbf{input\_collector}: Sistema de captura multi-dispositivo con arquitectura hexagonal, Core Event System framework-independiente, y soporte completo para teclado, ratón y controladores PS4
        \item \textbf{eeg\_processor}: Framework de procesamiento neuroseñal con arquitectura SOA, soporte para 20+ fabricantes de dispositivos EEG, análisis avanzado (ICA, wavelet, MNE-Python) y interfaces multi-modal
    \end{itemize}

    \item \textbf{Patrones arquitectónicos implementados}:
    \begin{itemize}
        \item \textbf{Event-Driven Architecture}: Sistema de eventos centralizado para comunicación desacoplada
        \item \textbf{Modelo-Vista-Controlador (MVC)}: Implementado en interfaces complejas con controladores especializados
        \item \textbf{Singleton Pattern}: Para servicios centralizados con thread-safety
        \item \textbf{Observer Pattern}: Para actualizaciones reactivas via Qt Signals/Slots
        \item \textbf{Strategy Pattern}: Para adaptación dinámica según contexto y dispositivos
        \item \textbf{Factory Pattern}: Para creación dinámica de configuraciones
    \end{itemize}

    \item \textbf{Integración inter-módulo avanzada}:
    \begin{itemize}
        \item \textbf{Lab Streaming Layer (LSL)}: Sincronización temporal de microsegundos
        \item \textbf{API REST centralizada}: Coordinación via api\_explorer
        \item \textbf{WebSocket streaming}: Comunicación tiempo real de alta frecuencia
        \item \textbf{Core Event Systems}: Comunicación framework-independiente
    \end{itemize}

    \item \textbf{Consideraciones de implementación críticas}:
    \begin{itemize}
        \item \textbf{Thread safety}: Estrategias granulares de sincronización
        \item \textbf{Optimizaciones de rendimiento}: Memory pooling, circular buffers, vectorización
        \item \textbf{Manejo de errores}: Circuit breaker, graceful degradation, error isolation
        \item \textbf{Escalabilidad}: Arquitectura preparada para extensiones futuras
    \end{itemize}
\end{enumerate}

El diseño resultante proporciona una base para el desarrollo de interfaces cerebro-computadora, cumpliendo con los objetivos de:

\begin{itemize}
    \item \textbf{Modularidad}: Componentes independientes con interfaces bien definidas
    \item \textbf{Extensibilidad}: Fácil adición de nuevos dispositivos, algoritmos y funcionalidades
    \item \textbf{Escalabilidad}: Arquitectura preparada para cargas de trabajo variables
    \item \textbf{Mantenibilidad}: Código bien estructurado con patrones establecidos
    \item \textbf{Usabilidad}: Múltiples interfaces adaptadas a diferentes tipos de usuarios
    \item \textbf{Rendimiento}: Optimizaciones específicas para datos de alta frecuencia
\end{itemize}

La arquitectura modular facilita el desarrollo incremental, la testabilidad individual de componentes, y el mantenimiento a largo plazo del sistema.
